%%
%% This is file `sigconf.tex',
%% generated with the docstrip utility.
%%
%% The original source files were:
%%
%% samples.dtx  (with options: `all,proceedings,sigconf')
%%
%% IMPORTANT NOTICE:
%%
%% For the copyright see the source file.
%%
%% Any modified versions of this file must be renamed
%% with new filenames distinct from sigconf.tex.
%%
%% For distribution of the original source see the terms
%% for copying and modification in the file samples.dtx.
%%
%% This generated file may be distributed as long as the
%% original source files, as listed above, are part of the
%% same distribution. (The sources need not necessarily be
%% in the same archive or directory.)
%%
%%
%% Commands for TeXCount
%TC:macro \cite [option:text,text]
%TC:macro \citep [option:text,text]
%TC:macro \citet [option:text,text]
%TC:envir table 0 1
%TC:envir table* 0 1
%TC:envir tabular [ignore] word
%TC:envir displaymath 0 word
%TC:envir math 0 word
%TC:envir comment 0 0
%%
%% The first command in your LaTeX source must be the \documentclass
%% command.
%%
%% For submission and review of your manuscript please change the
%% command to \documentclass[manuscript, screen, review]{acmart}.
%%
%% When submitting camera ready or to TAPS, please change the command
%% to \documentclass[sigconf]{acmart} or whichever template is required
%% for your publication.
%%
%%
\documentclass[sigconf]{acmart}
%%
%% \BibTeX command to typeset BibTeX logo in the docs
\AtBeginDocument{%
  \providecommand\BibTeX{{%
    Bib\TeX}}}

%% Rights management information.  This information is sent to you
%% when you complete the rights form.  These commands have SAMPLE
%% values in them; it is your responsibility as an author to replace
%% the commands and values with those provided to you when you
%% complete the rights form.
\setcopyright{acmlicensed}
\copyrightyear{2018}
\acmYear{2018}
\acmDOI{XXXXXXX.XXXXXXX}
%% These commands are for a PROCEEDINGS abstract or paper.
\acmConference[Conference acronym 'XX]{Make sure to enter the correct
  conference title from your rights confirmation email}{June 03--05,
  2018}{Woodstock, NY}
%%
%%  Uncomment \acmBooktitle if the title of the proceedings is different
%%  from ``Proceedings of ...''!
%%
%%\acmBooktitle{Woodstock '18: ACM Symposium on Neural Gaze Detection,
%%  June 03--05, 2018, Woodstock, NY}
\acmISBN{978-1-4503-XXXX-X/2018/06}











% add usepackage by yma391
\usepackage[table]{xcolor}
\definecolor{singleviewbg}{RGB}{255,243,224}   % light orange
\definecolor{multiviewbg}{RGB}{227,242,253}    % light blue
\definecolor{ynote}{RGB}{142,36,170}
\newcommand{\ynote}[1]{\textcolor{ynote}{#1}}  % zhuoxing's comments

\usepackage{tikz}
\usepackage{pgfplots}
\pgfplotsset{compat=1.18,
  tpchaxis/.style={
    width=\columnwidth,
    axis line style={black!60},
    xtick pos=bottom, ytick pos=left,
    tick align=outside, tick style={black!60},
    ymajorgrids, grid style={black!12, very thin},
    tick label style={font=\scriptsize},
    label style={font=\small},
  }}
\definecolor{tpchS}{RGB}{31,90,166}
\definecolor{tpchQ}{RGB}{204,85,0}
\definecolor{tpchT}{RGB}{16,133,120}
\definecolor{tpchR}{RGB}{176,32,66}
\usepackage{enumitem}
% Query Schema
\usepackage{booktabs}
% \usepackage{amssymb}
\usepackage{graphicx}
\graphicspath{{./}{tpch_database_experiment/}}
\newcommand{\rot}[1]{\rotatebox{90}{\textbf{#1}}}
\newcommand{\usemark}{\scalebox{1.25}{$\bullet$}}





%%
%% Submission ID.
%% Use this when submitting an article to a sponsored event. You'll
%% receive a unique submission ID from the organizers
%% of the event, and this ID should be used as the parameter to this command.
%%\acmSubmissionID{123-A56-BU3}

%%
%% For managing citations, it is recommended to use bibliography
%% files in BibTeX format.
%%
%% You can then either use BibTeX with the ACM-Reference-Format style,
%% or BibLaTeX with the acmnumeric or acmauthoryear sytles, that include
%% support for advanced citation of software artefact from the
%% biblatex-software package, also separately available on CTAN.
%%
%% Look at the sample-*-biblatex.tex files for templates showcasing
%% the biblatex styles.
%%
%%
%% The majority of ACM publications use numbered citations and
%% references.  The command \citestyle{authoryear} switches to the
%% "author year" style.
%%
%% If you are preparing content for an event
%% sponsored by ACM SIGGRAPH, you must use the "author year" style of
%% citations and references.
%% Uncommenting
%% the next command will enable that style.
%%\citestyle{acmauthoryear}


%%
%% end of the preamble, start of the body of the document source.
\begin{document}

%%
%% The "title" command has an optional parameter,
%% allowing the author to define a "short title" to be used in page headers.
\title{Normalization as a Safer Default:
A Neo4j TPC-H Case Study}

%%
%% The "author" command and its associated commands are used to define
%% the authors and their affiliations.
%% Of note is the shared affiliation of the first two authors, and the
%% "authornote" and "authornotemark" commands
%% used to denote shared contribution to the research.

%%
%% By default, the full list of authors will be used in the page
%% headers. Often, this list is too long, and will overlap
%% other information printed in the page headers. This command allows
%% the author to define a more concise list
%% of authors' names for this purpose.
%%
%% The abstract is a short summary of the work to be presented in the
%% article.
\begin{abstract}
Property graphs make denormalization attractive because copying properties
across relationship boundaries can replace traversals with local access.
However, each many-to-one fold studied here also induces a functional
dependency and replicates dependent values, so its update cost and query
benefit need not vary monotonically with the number of removed schema links.
We evaluate normalization as a conservative default in Neo4j at TPC-H
SF~0.01. Our first study constructs 59 role-aware, lossless replacement
layouts and compares each with a complete normalized eight-table layout under
the same 86,805-fact scalar-update workload. The normalized layout has lower
observed update time in 55 of 59 cases, including all eight direct folds and
47 of 51 recursive folds. The replacement-to-normalized time ratio ranges from
$0.965\times$ to $2.622\times$ (median $1.684\times$), while the replacement
requires $1.000230$--$6.938\times$ as many physical property-cell writes (median
$3.495\times$). A separate destructive conversion measures complete
replacement-to-normalized reconstruction rather than one-step unfolding.
Our second study retains the normalized graph and evaluates 74 query-specific
materialized-fold rewrite strategies over 16 join queries. In the complete-case
five-binding analysis, denormalization provides a query-execution benefit for
specific candidates: 24 of 69 complete strategies have a lower folded mean,
and Q3 and Q8 have lower folded macro-averages. The benefit is not systematic;
the normalized baseline has a lower query-level macro-average for 14 of 16
queries and a lower strategy mean for 45 of 69 complete strategies. The single
observed build times are on a subsecond-to-seconds scale while the query
differences are measured in milliseconds. These descriptive results support a
normalized source of truth augmented by candidate-specific denormalization
when its execution savings are validated and its construction and maintenance
costs can be amortized. The study does not measure materialized-view refresh
and is limited by its small scale, one measured update run per layout, one
measured conversion without a warm-up, and fixed query phase order.
\end{abstract}

%%
%% Keywords. The author(s) should pick words that accurately describe
%% the work being presented. Separate the keywords with commas.
\keywords{property graphs, normalization, denormalization, functional
dependencies, update amplification, materialized views, query rewriting,
Neo4j, TPC-H}

%%
%% This command processes the author and affiliation and title
%% information and builds the first part of the formatted document.
\maketitle


\section{Introduction}

Normalization aims to give each fact one authoritative representation and
connect related objects through explicit schema links. A property graph also
permits the opposite physical choice: properties from a referenced object can
be copied into referencing nodes, or a joined subgraph can be materialized as
a view. Such denormalization may remove traversals from a query, but it replicates
functionally dependent values across the join fan-out and introduces build,
storage, update, and maintenance obligations.

The design question is therefore not whether a fold can remove links, but
whether that fold produces a repeatable net benefit for the workload that will
use and maintain it. We ask whether denormalization consistently improves
property-graph performance, or whether the normalized graph should remain the
default and folds should be admitted only as validated, workload-specific
optimizations. In this paper, a ``safer default'' means the design with lower
observed exposure to logical value duplication, replicated writes, and
unprofitable query rewrites under our protocol; it is not a claim of universal
or statistically established superiority.

We operationalize a conceptual fold framework in two complementary ways. A
lossless replacement layout exposes the footprint and full-database update
cost of crossing normalization boundaries. A coexisting materialized fold view
tests whether removing selected links from a query plan improves execution
enough to justify constructing a derived structure. Across the tested TPC-H
workloads, the complete normalized layout has lower observed full-workload
update time in 55 of 59 comparisons, while the
normalized query has a lower observed mean for most query groups and most
complete rewrite strategies.
On the query side, however, 24 complete strategies have lower folded means, and
Q3 and Q8 have lower folded macro-averages in the complete-case analysis. These
positive cases show that a suitably matched fold view can reduce execution
time. The evidence therefore supports a
normalized core with selective, workload-specific denormalization rather than
either universal normalization or blanket denormalization.

This study makes three contributions:
\begin{itemize}[leftmargin=*]
  \item It maps conceptual folds and their induced functional dependencies to
  lossless replacement property-graph layouts and coexisting materialized Join
  Views.
  \item It evaluates 59 direct and recursive folded layouts for logical
  footprint, full-database update amplification, and complete
  replacement-to-normalized conversion, and 74 query-specific fold-view
  rewrites over 16 TPC-H queries.
  \item It motivates a conservative, scope-bound decision rule: retain the
  normalized graph as the default source of truth, and require each
  denormalized candidate to demonstrate semantic equivalence, repeatable
  workload benefit, and cost amortization before adoption.
\end{itemize}


\section{Experimental Evaluation}

We treat the normalized graph as the reference design and evaluate each fold
as a departure that must justify its redundancy and operational cost. The
study asks three questions. First, how do role-aware folds and join fan-out
translate into logical footprint and full-database update amplification?
Second, how many repeated update workloads are required to repay the cost of
complete replacement-to-normalized conversion?
Third, when do materialized-fold rewrites outperform their paired normalized
queries, and are those execution-time savings large enough to amortize view
construction cost?

The update experiment compares two complete, run-scoped representations: a
normalized clone of all eight TPC-H relations and a lossless replacement layout
whose selected relations are represented by the fold, residual rows, and any
required companion structures. It also times a destructive conversion from
the replacement layout to the complete normalized layout. The query experiment
instead retains the normalized base graph and adds a candidate materialized
Join View. Together, they expose complementary costs relevant to adoption, but
they do not form a combined end-to-end cost test. The replacement-layout update
is not a measurement of materialized-view refresh, which is outside this suite.

\subsection{Experimental Environment}
\label{subsec:fd_experimental_environment}

Both studies use TPC-H SF~0.01; the full-layout update and normalization study
uses the \texttt{tpch-sf-001} database. All generated
nodes and relationships are run-scoped, and the imported TPC-H graphs remain
read-only.
The redundancy report records Python~3.13.14, Neo4j Python driver~6.2.0, and
Neo4j Enterprise~2026.06.0, accessed locally through Bolt~6.0 with the
server-default Cypher runtime; the query run file does not retain corresponding
client/server version metadata. Reported operation and query times are client
wall-clock times through result consumption and commit.

\subsection{Experimental Methodology}
\label{subsec:fd_redundancy_method}

\paragraph{Conceptual-to-physical mapping.}
An E/R link from a higher-order object $O$ to a lower-order object $O'$
separates the two object types and thereby eliminates the redundancy that
$O'$ would introduce if it were represented inside $O$. Following the fold
definition, folding $O'$ into $O$ removes this link and induces on the
resulting object $O_f$ the functional dependency
$\textit{id}(O')\rightarrow
\textit{comp}(O')\cup\textit{attr}(O')$.
Our property-graph implementation is a projected physical realization of this
definition. Property graphs represent inter-object connections explicitly as
relationships, allowing the exposed topology of a fold to be preserved by
reconnecting its boundary relationships. Within a folded node, the existing
foreign-key properties of the referencing object already identify the
referenced object and therefore the referenced primary key is not copied
again. Only the referenced non-key properties are added to the folded node.

The measured duplication consequently concerns the dependent non-key
properties determined by $\textit{id}(O')$. For the many-to-one TPC-H folds
evaluated here, $\textit{id}(O')$ is not a key of $O_f$; the same referenced
value is therefore stored once for every root-grain node in that key's
fan-out. Recursive folds repeat this operation as additional object types are
incorporated. Schema-level fold degree organizes the conceptual design space,
but the realized redundancy is determined by instance-level fan-out and
property width.


\paragraph{Candidate folds.}
The normalized TPC-H schema contains eight schema-defined E/R link types, and
hence admits $2^8$ schema-level link subsets conceptually. The operational
suite neither enumerates nor forms a literal subset of that Boolean lattice.
It expands shared object types into branch-specific occurrences and evaluates
the eight one-occurrence foreign-key--primary-key folds plus 51 connected,
rooted constructions in that role-expanded occurrence space. Strategy
initials denote TPC-H relations; N1/R1 and N2/R2 distinguish the
\textsc{Nation}/\textsc{Region} occurrences on the customer and supplier
branches, respectively. Applying the shared \textsc{Nation}--\textsc{Region}
link separately in these two branches permits a construction with nine fold
occurrences despite there being only eight schema link types. We therefore use
$d$ for operational fold-chain length (the number of fold occurrences), not
for the number of distinct schema-link types; the 51 recursive constructions
have $d=2$--9. Every stage preserves the root cardinality and carries forward
only the boundary relationships required by a later declared fold.

\paragraph{Lossless replacement, full-database updates, and normalization.}
Let $G_N$ be a run-scoped, fully normalized clone of all eight TPC-H relations
and their eight foreign-key relationship types. Let $G_R$ be the corresponding
lossless replacement layout. In $G_R$, the selected relations are represented
by the folded primary component; source rows that do not participate in the
fold are retained as residual components, shared \textsc{Nation}/\textsc{Region}
identities are represented by the required companion or role-scoped
occurrences, and the remaining relations and complete external topology are
preserved. Both layouts reconstruct the same eight physical relations. N1/N2
and R1/R2 are storage occurrences of the same physical \textsc{Nation} and
\textsc{Region} facts, not additional logical facts. The imported source graph
is used to construct and validate the isolated layouts but remains read-only.

Table~\ref{tab:fd_normalization_results} reports the actually materialized
replacement-minus-normalized changes
\[
  \Delta V=|V(G_R)|-|V(G_N)|,\qquad
  \Delta E=|E(G_R)|-|E(G_N)|,
\]
and
\[
  \Delta C=C(G_R)-C(G_N),
\]
where $C(G)$ counts logical TPC-H business-property cells, including declared
primary- and foreign-key columns but excluding relationship properties and
experimental bookkeeping. At SF~0.01, $G_N$ contains 86,805 nodes, 152,975
relationships, and 1,168,615 logical property cells in every case. These
metrics describe complete, lossless layouts actually used by the full-update
experiment; they are not the incremental footprint of a coexisting view or
Neo4j store bytes.

The shared scalar workload contains one logical fact for every row of every
physical TPC-H relation, for 86,805 facts in total. It cyclically changes
\texttt{r\_name}, \texttt{n\_name}, \texttt{s\_phone},
\texttt{c\_mktsegment}, \texttt{p\_name}, \texttt{ps\_supplycost},
\texttt{o\_orderpriority}, and \texttt{l\_shipmode}. A mapping is generated
once from each complete source relation's active domain and then applied to
both layouts. In $G_N$, each logical fact has one physical property-cell
target; in $G_R$, the same fact is propagated to every primary, residual,
companion, or unaffected occurrence that represents it. Thus the logical work
is identical while the physical property-write count is allowed to differ.
Foreign-key and topology updates are outside this scalar workload.

Each layout receives one warm-up and one measured update run. For paired run
$i$, let $t_{N,i}$ and $t_{R,i}$ be the complete eight-table client wall times.
The table reports arithmetic means
\[
 t_N=\operatorname{mean}_i(t_{N,i}),\qquad
 t_R=\operatorname{mean}_i(t_{R,i}),
\]
the paired mean difference
\[
 \Delta_U=\operatorname{mean}_i(t_{N,i}-t_{R,i}),
\]
and the paired mean time ratio
\[
 A_T=\operatorname{mean}_i(t_{R,i}/t_{N,i}).
\]
Accordingly, $\Delta_U>0$ favors replacement, whereas $A_T>1$ favors
normalization. With one measured run in the current report, each mean equals
its single observation. Layout construction, routing preparation, restoration,
and validation are excluded from update timing; caches are not cleared.

After both update arms have been restored and validated, a separate,
destructive conversion reconstructs the eight normalized relations from
$G_R$, rebuilds all eight foreign-key relationship types from the converted
foreign-key properties, and deletes the replacement input. Its client wall
time, denoted $t_{R\rightarrow N}$ and labelled ``Normalization'' in the table,
includes result consumption and commits. The timed conversion does not read
the source graph. Removal of the prebuilt normalized comparison layout and
routing anchors, index DDL, post-conversion validation, source fingerprinting,
and final cleanup are excluded. The conversion is measured once without a
warm-up, with warm caches and pre-existing lookup indexes.

When normalization saves update time, i.e.,
$\overline{s}=\operatorname{mean}_i(t_{R,i}-t_{N,i})>0$, we report the
analyzer-derived point estimate
\[
  \kappa=\left\lceil
  \frac{t_{R\rightarrow N}}{\overline{s}}
  \right\rceil .
\]
It is the minimum number of subsequent complete eight-table update workloads
whose observed mean saving would repay the measured conversion. It is N/A
when $\overline{s}\leq0$. This conversion threshold is distinct from a
materialized view's query/build or query/refresh recoup ratio.

\paragraph{Materialized-fold query rewrites.}
For 16 TPC-H join queries with at least one schema-driven candidate, we test 74
query-specific fold-view rewrite strategies generated from connected subgraphs
of the query's schema-link footprint. One flattened view node is stored per
output row of the candidate's canonical join tree, with the boundary
relationships needed by the remainder of the query. A full-footprint
candidate supports a projection or aggregation over the view; a partial
candidate replaces its covered links and leaves a reduced join against the
base graph.

Each view is built and removed one strategy at a time. Its normalized query
and folded-view rewrite are executed with the same five generated parameter
bindings, with the normalized phase preceding the rewrite phase. No separate
timing warm-up is discarded, and caches are not explicitly cleared. For every
uncensored, completed pair, the complete result bags are compared while
preserving duplicate and null multiplicities and using tolerances only for
floating-point values. Each rewrite is subject to an approximate server-side
timeout set to twice its paired normalized runtime.
Runs that time out are right-censored, whereas executions that finish after the
nominal limit retain their observed times; the configured cap is recorded as a
runtime lower bound for each censored run. A strategy with any missing or
censored timing is classified as incomplete and excluded from the aggregate.
For every complete strategy, normalized and folded client times are each
averaged across the five distinct parameter bindings; the resulting strategy
means then receive equal weight in the query-level macro-average.
No fastest-view selection or baseline fallback is applied. We define
$\Delta_Q=t_N-t_F$ and $S_Q=t_N/t_F$, where $t_N$ and $t_F$ are normalized and
folded-rewrite times. Hence, $\Delta_Q>0$ and $S_Q>1$ indicate that the view
rewrite is faster. View build time is measured separately and excluded from
query time; each strategy contributes one measured build time.
Materialized-view refresh and maintenance are not measured in this suite.
For each complete strategy $s$, let $\bar t_{N,s}$ and $\bar t_{F,s}$ denote
its normalized and folded means over the five parameter bindings, and let $b_s$
be its single measured view-build time. When the candidate-specific saving
$g_s=\bar t_{N,s}-\bar t_{F,s}$ is positive, we define
\[
  \widetilde{\kappa}_{B,s}=\frac{b_s}{g_s},
  \qquad
  \kappa_{B,s}=\left\lceil\widetilde{\kappa}_{B,s}\right\rceil .
\]
Build and query times use the same unit in this calculation. The threshold
assumes that future executions reproduce the observed five-binding mean saving
$g_s$; it ignores refresh, storage, and timing variability.
Accordingly, $g_s>0$ establishes that candidate $s$ has a lower observed folded
mean over the five tested bindings, but it is not by itself sufficient for
adoption: the rewrite must also preserve the result bag, improve representative
executions under repeated measurement, and repay its own construction and
maintenance costs.
Because refresh is unmeasured, this study identifies query-time opportunities
rather than end-to-end profitability.

\subsection{Results and Discussion}
\label{subsec:fd_redundancy_results}

\subsubsection{Full-layout update and normalization}

\begin{table}[t]
\centering
\scriptsize
\setlength{\tabcolsep}{2.8pt}
\caption{Full eight-table update, replacement-to-normalized conversion, and
actual lossless-layout changes for the 26 folds with chain length $d\leq3$ at
TPC-H SF~0.01.}
\label{tab:fd_normalization_results}
\resizebox{0.48\textwidth}{!}{%
\begin{tabular}{rlrrrrrrrrr}
\toprule
No. & Fold & $t_N$ & $t_R$ & $\Delta_U$ & $A_T$ &
Norm. & $\kappa$ & $\Delta V$ & $\Delta E$ & $\Delta C$ \\
\midrule
 1 & N-R        & 1,216 & 1,219 &    -2 & 1.002 &   5,467 & 2,439 &     -5 &    -25 &   +35 \\
 2 & S-N        & 1,131 & 1,142 &   -11 & 1.010 &   3,682 &   341 &      0 &      0 &  +300 \\
 3 & C-N        & 1,145 & 1,240 &   -94 & 1.082 &   3,638 &    39 &      0 &      0 &   +5K \\
 4 & O-C        & 1,148 & 1,201 &   -53 & 1.046 &   4,260 &    81 &    -1K &    -1K &  +97K \\
 5 & PS-P       & 1,116 & 1,229 &  -113 & 1.101 &   3,504 &    32 &    -2K &    -8K &  +46K \\
 6 & PS-S       & 1,147 & 1,148 &    -1 & 1.001 &   3,516 & 5,521 &   -100 &   -100 &  +47K \\
 7 & L-O        & 1,127 & 1,279 &  -152 & 1.135 &   3,756 &    25 &   -15K &   -15K & +346K \\
 8 & L-PS       & 1,185 & 1,322 &  -138 & 1.116 &   4,877 &    36 &    -8K &   +44K & +141K \\
 9 & L-O-C      & 1,154 & 1,598 &  -444 & 1.385 &   4,456 &    11 &   -16K &   -16K & +760K \\
10 & L-O-PS     & 1,199 & 1,591 &  -392 & 1.327 &   5,179 &    14 &   -23K &   +29K & +487K \\
11 & L-PS-S     & 1,262 & 1,723 &  -461 & 1.365 &   4,916 &    11 &    -8K &   +44K & +501K \\
12 & L-PS-P     & 1,328 & 1,654 &  -327 & 1.246 & 132,296 &   406 &   -10K &   -16K & +604K \\
13 & O-C-N1     & 1,332 & 1,305 &   +26 & 0.980 &   4,026 &   N/A &    -1K &    -1K & +142K \\
14 & C-N1-R1    & 1,257 & 1,213 &   +43 & 0.965 &   3,818 &   N/A &     -5 &    -2K &   +8K \\
15 & PS-S-P     & 1,269 & 1,246 &   +23 & 0.982 &   3,554 &   N/A &    -2K &    -8K &  +93K \\
16 & PS-S-N2    & 1,229 & 1,341 &  -111 & 1.091 &   3,755 &    34 &   -100 &   -100 &  +71K \\
17 & S-N2-R2    & 1,205 & 1,177 &   +27 & 0.977 &   3,846 &   N/A &     -5 &   -125 &  +535 \\
18 & L-O-C-PS   & 1,220 & 1,802 &  -582 & 1.477 &   5,119 &     9 &   -24K &   +28K & +900K \\
19 & L-O-C-N1   & 1,224 & 1,788 &  -565 & 1.461 &   4,742 &     9 &   -16K &   -16K & +940K \\
20 & L-O-PS-S   & 1,233 & 1,809 &  -575 & 1.467 &   4,650 &     9 &   -23K &   +29K & +847K \\
21 & L-O-PS-P   & 1,169 & 1,801 &  -632 & 1.540 &   4,527 &     8 &   -25K &   -31K & +950K \\
22 & L-PS-S-P   & 1,186 & 1,796 &  -610 & 1.514 &   4,743 &     8 &   -10K &   -16K & +964K \\
23 & L-PS-S-N2  & 1,211 & 1,786 &  -575 & 1.475 &   5,181 &    10 &    -8K &   +44K & +681K \\
24 & O-C-N1-R1  & 1,208 & 1,335 &  -127 & 1.106 &   3,978 &    32 &    -1K &   -16K & +172K \\
25 & PS-S-P-N2  & 1,205 & 1,304 &   -98 & 1.082 &   3,774 &    39 &    -2K &    -8K & +117K \\
26 & PS-S-N2-R2 & 1,161 & 1,181 &   -20 & 1.017 &   3,657 &   184 &   -105 &    -8K &  +87K \\
\bottomrule
\end{tabular}
}
\textit{Note:} All time columns report client wall-clock milliseconds.
$t_N$, $t_R$, $\Delta_U$, and $A_T$ are paired-run arithmetic means as defined
in the text; the current report has one measured update run per layout. The
four time columns are independently rounded to the nearest
millisecond from their unrounded means; consequently, subtracting the displayed
$t_N$ and $t_R$ can differ from the displayed $\Delta_U$ by 1\,ms. Positive
$\Delta_U$ favors replacement, whereas $A_T>1$ favors normalization.
$t_{R\rightarrow N}$ is the separately measured, destructive
full replacement-to-normalized client time. $\kappa$ is the minimum number of
subsequent complete eight-table update workloads needed to repay that
conversion under the observed mean saving; N/A denotes a non-positive
normalization saving. $\Delta V$, $\Delta E$, and $\Delta C$ are actual
replacement-minus-normalized changes in nodes, explicit relationships, and
logical business-property cells. K denotes $10^3$ and M denotes $10^6$;
suffixed changes are rounded to the nearest K or M. They are logical counts,
not Neo4j store bytes or the incremental footprint of a coexisting view.
\end{table}

Table~\ref{tab:fd_normalization_results} shows the 26 configurations of chain
length at most three; the aggregate statistics below cover all 59 cases. The
complete normalized layout has lower observed full-workload update time in 55
cases: all eight direct folds and 47 of the 51 recursive folds. The four
replacement-faster cases are O--C--N1, C--N1--R1, PS--S--P, and S--N2--R2.
Across all cases, $A_T$ ranges from $0.965$ to $2.622$, with a median of
$1.684$. The paired mean difference $\Delta_U$ ranges from
$-2{,}294.966$ to $+43.490$\,ms (median $-843.131$\,ms); hence the largest
observed differences favor normalization, whereas the four replacement
advantages are comparatively small. Since each arm currently has only one
measured run, these are descriptive observations rather than estimates of
run-to-run stability.

The physical write counts provide the structural mechanism behind this
ordering. The normalized arm always writes 86,805 property cells, one per
logical fact. Across the 59 replacement layouts, the physical property-write
ratio ranges from $1.000230$ to $6.938$ (median $3.495$). Unlike the former
final-layer $J/D$ diagnostic, this ratio covers the entire eight-table workload
and all occurrences in the lossless replacement. More physical targets need
not produce a proportionate wall-clock increase because table sizes, index
lookups, transaction boundaries, caching, and orchestration also contribute to
the measured time. Nevertheless, the normalized arm is the more robust update
choice in the observed cases.

The actual lossless-layout footprint also separates graph-object counts from
copied logical state. $\Delta V$ is negative in 57 cases and zero in two
(range $-26{,}126$ to 0; median $-23{,}096$). $\Delta E$ is negative in 45
cases, zero in two, and positive in 12 (range $-152{,}467$ to $+44{,}183$;
median $-16{,}025$). By contrast, $\Delta C$ is positive in all 59 cases,
ranging from 35 to $2{,}325{,}555$ logical property cells with a median of
$1{,}060{,}535$. Relative to the fixed normalized baseline of 1,168,615 cells,
the replacement cell ratio ranges from 1.000030 to 2.990010 (median 1.907514).
Thus a fold can reduce node and relationship counts while still increasing the
amount of logical property state that must be represented.

Full replacement-to-normalized conversion ranges from 3,504.051 to
133,026.518\,ms, with a median of 5,389.725\,ms. The maximum is a pronounced
single-run outlier, so the median is more representative of the observed case
distribution than the mean. Normalization saves update time in 55 cases; for
these cases, the analyzer-derived integer threshold $\kappa$ ranges from 4 to
5,521 complete update workloads, with a median of 8. Thirty-three thresholds
are at most ten workloads and 37 are at most fifteen. The threshold is N/A for
the four replacement-faster cases because conversion cannot be repaid by a
non-positive subsequent normalization saving. These point estimates assume
repetition of exactly the same complete scalar workload and timing boundary.
Moreover, conversion is measured once without a warm-up and with indexes
already online; its thresholds should therefore not be interpreted as stable
deployment break-even guarantees.

\subsubsection{Conditional query benefits from materialized folds}

\begin{table}[t]
\centering
\scriptsize
\caption{Normalized-baseline and materialized-fold query performance for
complete strategies.}
\label{tab:query_rewrite_performance}
\resizebox{0.47\textwidth}{!}{
\begin{tabular}{c c r r r r r}
\toprule
Q   & C/T & F-Que. & N-Que. & $\Delta_Q$ & $S_Q$ & MV Build \\
\midrule
Q2  & 8/10  & 70.1  & 26.2  & $-43.9$  & 0.37 & 1.54 \\
Q3  & 2/3   & 86.7  & 123.7 & $+37.0$  & 1.43 & 9.15 \\
Q4  & 1/1   & 108.8 & 48.8  & $-60.1$  & 0.45 & 6.24 \\
Q5  & 6/6   & 80.0  & 57.1  & $-22.9$  & 0.71 & 6.83 \\
Q7  & 7/7   & 48.8  & 43.4  & $-5.4$   & 0.89 & 9.94 \\
Q8  & 11/11 & 84.1  & 102.6 & $+18.5$  & 1.22 & 9.84 \\
Q9  & 4/4   & 67.7  & 58.1  & $-9.6$   & 0.86 & 11.55 \\
Q10 & 6/6   & 107.1 & 51.7  & $-55.4$  & 0.48 & 9.28 \\
Q11 & 2/3   & 10.4  & 8.1   & $-2.3$   & 0.77 & 0.91 \\
Q12 & 1/1   & 146.9 & 122.1 & $-24.9$  & 0.83 & 6.11 \\
Q13 & 1/1   & 125.6 & 115.0 & $-10.5$  & 0.92 & 1.42 \\
Q16 & 3/3   & 90.6  & 58.3  & $-32.2$  & 0.64 & 1.43 \\
Q18 & 3/3   & 370.8 & 178.4 & $-192.4$ & 0.48 & 8.56 \\
Q20 & 12/12 & 34.6  & 21.1  & $-13.5$  & 0.61 & 10.84 \\
Q21 & 1/2   & 54.7  & 52.6  & $-2.0$   & 0.96 & 0.49 \\
Q22 & 1/1   & 45.0  & 23.7  & $-21.4$  & 0.53 & 1.67 \\
\bottomrule
\end{tabular}
}
\textit{Note:} C/T denotes complete/tested strategies. F-Que. and N-Que. are
folded-rewrite and normalized query times in milliseconds; MV Build is in
seconds and is the mean over the same complete strategies. $\Delta_Q=t_N-t_F$
and $S_Q=t_N/t_F$, so positive $\Delta_Q$ and $S_Q>1$ favor the fold view.
Negative $\Delta_Q$ and $S_Q<1$ favor normalization. Differences and ratios
are computed from unrounded macro-means.
\end{table}

The query study gives a qualified affirmative answer: denormalization can
reduce query-execution time when a fold view is matched to a suitable query.
Of the 74 tested strategies, 69 complete all five paired executions, and 24 of
those 69 have a lower folded five-binding mean. At the query-group level, Q3
and Q8 have folded macro-averages lower than their normalized baselines by
37.0\,ms ($1.43\times$) and 18.5\,ms ($1.22\times$), respectively. These are
observed execution-time opportunities under the tested bindings. The 24
lower-mean fold candidates span nine of the 16 query groups, including seven
groups whose candidate macro-average still favors normalization. Thus, a
normalized-favoring macro-average does not imply that a query has no useful
fold candidate.

The benefit is selective rather than general. The normalized five-binding mean
is lower for the other 45 complete strategies, and the normalized baseline has
the lower query-level macro-average for 14 of the 16 query groups. The
normalized query is therefore the lower-risk fallback when candidate-specific
evidence is unavailable; a fold should be chosen because its own measurements
show a benefit, not merely because it removes schema links.

The Q2 folded macro-average is 2.67 times the normalized macro-average
($S_Q=0.37$), the largest relative ratio, while Q18 has the largest absolute
difference at 192.4\,ms. Seven capped runs make five strategies
incomplete. Omitting them makes
Table~\ref{tab:query_rewrite_performance} conditional on completion and
optimistic toward folded rewrites. As a performance-only sensitivity check,
substituting each capped run's recorded runtime lower bound makes every omitted
strategy's folded five-binding lower bound higher than its normalized mean.
This yields at least 50 normalized-favoring runtime comparisons among the 74
tested strategies. When the lower-bound means are included in the query-level
macros, Q3 changes sign ($t_N=99.7$\,ms versus
$t_F^{\mathrm{LB}}=113.7$\,ms), so 15 of 16 favor normalized execution and
only Q8 retains a folded macro-average advantage. The capped executions lack a
complete validated result bag and therefore remain excluded from the main
aggregate.
The fixed phase order and absence of a discarded warm-up further make the
reported timings descriptive rather than causal estimates.

A fold can remove schema links from a query plan without thereby guaranteeing
lower physical runtime. Root grain, predicate selectivity, aggregation,
remaining traversals, cache state, and optimizer choices are plausible factors
that may determine whether the saved links translate into a benefit; this
experiment does not isolate their individual effects. Nor does an
execution-time benefit automatically imply a lower first-use cost. Among the
24 complete strategies with positive query-time savings in their five-binding
means, the candidate-specific build-only threshold $\kappa_{B,s}$ ranges from
3 to 4,022 executions, with a median of 885. Under this five-binding mean model,
$b_s+\bar t_{F,s}>\bar t_{N,s}$ for all 69 complete strategies: none repays its
measured build in one execution charged at the strategy's mean saving. A
positive candidate can nevertheless repay construction under sufficiently
repeated reuse, and this mean-based calculation does not rule out a favorable
individual binding. The lower endpoint illustrates the potential: the Q2
Nation--Region candidate has $\bar t_{N,s}=86.7$\,ms,
$\bar t_{F,s}=25.9$\,ms, $b_s=153.4$\,ms, and $\kappa_{B,s}=3$, even though
Q2's cross-candidate macro-average favors normalization. The median threshold
of 885 shows that such rapid payback is not typical of the 24 positive-mean
candidates. Any required refresh would add an unmeasured recurring cost.
These thresholds pair each candidate's own build time and query saving; they
cannot be obtained from the cross-candidate macro-averages in
Table~\ref{tab:query_rewrite_performance}.

\subsection{Overall Discussion: Normalization as the Safer Default}

The results support normalization as the safer default under the tested
workloads, not as a universally faster representation. Here, ``safer'' has
three concrete meanings. A dependent value has one authoritative source,
direct update work does not scale with the number of folded copies, and every
query retains a validated normalized fallback without first constructing a
derived view. The update study demonstrates the first asymmetry directly:
the lossless replacement layouts require between $1.000230$ and $6.938\times$ as
many physical property-cell writes as the normalized layout for the same
86,805 logical facts, and 55 of 59 have higher observed full-workload update
time. All eight direct folds favor normalization; four recursive layouts favor
replacement in the single measured run. Fold-chain length organizes the
candidate space, but occurrence multiplicity, target routing, and system state
determine much of the operational exposure.

The query study establishes a complementary result: denormalization is a real
but selective read optimization. Twenty-four complete strategies have lower
folded five-binding means, and Q3 and Q8 have lower folded macro-averages in
the complete-case analysis. These positive cases show that materializing a
suitable fold can reduce execution time for a selected query; repeated reuse is
what may allow that recurring saving to amortize construction. Q8 retains its
query-level advantage under the censoring lower-bound sensitivity analysis,
whereas Q3's macro-average advantage is conditional on the complete-case set.
Most complete strategy means are nevertheless higher than the normalized
baseline, query benefit is not guaranteed by schema coverage, and the results
provide no monotonic guarantee with fold-chain length. The lower-risk
architecture is therefore a normalized source of truth augmented with
independently validated materialized folds for stable, frequently repeated
queries.

No end-to-end profitability claim follows from combining the two suites. The
update experiment studies complete lossless replacement layouts, whereas the query
experiment studies coexisting materialized views whose refresh cost is
unmeasured. The macro-averages and candidate counts are equal-weight summaries,
not a workload-frequency-weighted objective. We therefore adopt a
normalized-core policy: a fold view should be admitted only after demonstrating
(i) bag-equivalent answers, (ii) positive and repeatable candidate-specific
query savings, (iii) construction-cost amortization at the expected query
frequency, and (iv) acceptable storage and refresh cost. This study checks the
first condition on the completed tested
bindings and measures parts of the second and third, but not the fourth.
Without all four forms of evidence, the normalized query remains the defensible
fallback.

\section{Conclusion}

Within the tested Neo4j TPC-H workload, normalization is the safer default
design. Every tested lossless replacement layout increases the physical
property-cell writes for the shared full-database workload, and 55 of 59 have
higher observed update time than the complete normalized layout in the single
measured run. Complete replacement-to-normalized conversion has a median
observed client time of 5,389.725\,ms, and the median analyzer-derived payback
threshold among the 55 normalization-faster cases is eight repeated update
workloads. Query execution, however, can benefit from selective
denormalization: 24 of 69 complete fold-view strategies, spanning nine query
groups, had lower folded five-binding means, although the normalized mean was
lower for the other 45 strategies. At the macro level, Q3 and Q8 favor folding
in the complete-case analysis, whereas 14 of 16 query groups favor
normalization; only Q8 retains that macro-level fold advantage in the
lower-bound sensitivity analysis. Candidate-specific build-only thresholds
range from 3 to 4,022 mean-equivalent executions, so selected candidates may
amortize construction under frequent reuse even though none does so in one
execution charged at its five-binding mean saving.

The resulting design recommendation is not ``never denormalize,'' but
``denormalize selectively.'' Retaining the normalized graph localizes
authoritative updates and preserves a query fallback, while a carefully matched
materialized fold can accelerate a stable, repeated query. Such gains must be
validated across representative parameters and repeated measurements. Expected
reuse must be high enough to amortize construction and any required refresh,
while storage and update costs must remain acceptable.

The conclusion is deliberately scoped to the tested data, system, and
protocol. The full-layout update arms each have one measured run after one
warm-up, while the destructive conversion is measured once without a warm-up
under warm-cache, index-assisted conditions. Larger scale factors, repeated
and counterbalanced measurements, physical storage accounting,
materialized-view refresh, and workload-level view selection are required
before making broader performance claims.



%%
%% The next two lines define the bibliography style to be used, and
%% the bibliography file.
\bibliographystyle{ACM-Reference-Format}
\bibliography{sample-base}



\end{document}
\endinput
%%
%% End of file `sigconf.tex'.
